\begin{problem}{}{}

    Demostrar que si un conjunto $X$ tiene mínimo, este es único.
    Dicho más formalmente:
    demostrar que si existen $c$, $c' \in X$ tal que $c \leq x$ y $c' \leq x$ para todo $x \in X$,
    entonces $c = c'$.

\end{problem}
\vspace{1ex}

Supongamos que existen $c, c' \in X$ satisfaciendo las condiciones dadas.
Queremos ver que son iguales.
Para lograr esto, un camino inteligente es tratar de utilizar la propiedad antisimétrica del orden ($a \le b \land b \le a \implies a = b$), siempre y cuando podamos relacionar $c$ con $c'$.
Efectivamente, como ambos son mínimos del mismo conjunto, podemos razonar como sigue:

\begin{itemize}
    \item Como $c$ es mínimo, tenemos $c \le x$ para todo $x \in X$. Como $c' \in X$, en particular, concluímos que $c \le c'$.
    \item De manera análoga al argumento anterior, podemos concluir que $c' \le c$.
\end{itemize}

Juntando ambas conclusionem, tenemos que $c \le c' \land c' \le c$.
Con esto, la propiedad antisimétrica nos asegura que $c = c'$, como queríamos probar.
Así, concluímos la prueba.
